\documentclass[11pt,a4paper]{article}

\usepackage[margin=2.5cm]{geometry}
\usepackage{enumitem}
\usepackage{booktabs}
\usepackage{array}
\usepackage{xcolor}
\usepackage{titlesec}
\usepackage{parskip}
\usepackage{hyperref}
\usepackage{fancyhdr}
\usepackage{amsmath}
\usepackage{amssymb}

\definecolor{accent}{RGB}{30,90,180}
\definecolor{warn}{RGB}{180,60,30}
\definecolor{muted}{RGB}{100,100,100}

\hypersetup{colorlinks=true, linkcolor=accent, urlcolor=accent, breaklinks=true}
\def\UrlBreaks{\do\/\do-\do\_\do.}

\titleformat{\section}{\large\bfseries\color{accent}}{}{0em}{}[\titlerule]
\titleformat{\subsection}{\normalsize\bfseries}{}{0em}{}

\pagestyle{fancy}
\fancyhf{}
\rhead{\textcolor{muted}{\small SparseRenderFormer}}
\lhead{\textcolor{muted}{\small Team Action Steps}}
\cfoot{\thepage}

\newcommand{\role}[1]{\textbf{[#1]}}
\newcommand{\critical}[1]{\textcolor{warn}{\textbf{#1}}}
\newcommand{\sq}{\item[$\square$]}

\begin{document}

\begin{center}
  {\LARGE\bfseries SparseRenderFormer}\\[4pt]
  {\large Team Action Steps}\\[6pt]
  {\color{muted} Team size: 4\quad Transformer experience: beginner\quad Target: 11~weeks}
\end{center}

\vspace{6pt}
\noindent\critical{Critical path:} Start data rendering (Phase~2) in Week~2, not after
implementation. Blender renders are compute-bound and must run overnight on the GPU
throughout development. Waiting for code to finish before generating data costs two weeks.

\vspace{6pt}

%-----------------------------------------------------------------------
\section{Phase~0\quad Learning Sprint \textnormal{(Week~1, all~4)}}

One focused week. Everyone works through the list independently, then the team
discusses together.

\begin{itemize}[leftmargin=1.8em, itemsep=2pt]
  \sq Watch Andrej Karpathy ``Let's build GPT from scratch'' (YouTube, $\approx$2~hrs)
  \sq Read ``Attention Is All You Need'' --- Vaswani et al., NeurIPS 2017
  \sq Read RenderFormer paper --- arXiv:2505.21925 (SIGGRAPH 2025)
  \sq Read NSA paper --- arXiv:2502.11089 (direct design inspiration for GP-Sparse)
  \sq Read \texttt{about\_sparserenderformer.tex} in this repo
  \sq Read \texttt{CLAUDE.md} for method design and token format
  \sq Clone RenderFormer: \texttt{git clone https://github.com/microsoft/renderformer}
  \sq Run RenderFormer inference on the Cornell box scene
\end{itemize}

\noindent\textbf{Exit criterion:} every team member can explain what the coarse, selected,
and local branches do and why they reduce attention complexity from $O(N^2)$ to $O(N)$.

%-----------------------------------------------------------------------
\section{Phase~1\quad Engineering Setup \textnormal{(Weeks~1--2, split by role)}}

\subsection{\role{P1} ML infrastructure}
\begin{itemize}[leftmargin=1.8em, itemsep=2pt]
  \sq Confirm GPU allocation (A100-80\,GB recommended; A6000 workable for ablations)
  \sq Set up Python environment: PyTorch~2.x, Triton $\geq$~3.x
  \sq Install Weights \& Biases (or MLflow) for experiment tracking
  \sq Back up RenderFormer pretrained checkpoint to shared storage before fine-tuning
\end{itemize}

\subsection{\role{P2} Baseline reproduction}
\begin{itemize}[leftmargin=1.8em, itemsep=2pt]
  \sq Download checkpoint \texttt{microsoft/renderformer-v1.1-swin-large} from HuggingFace
  \sq Run inference on one Objaverse scene end to end
  \sq Reproduce one PSNR number from Table~1 of the RenderFormer paper
  \sq Document the exact command used and the result in a shared note
\end{itemize}

\subsection{\role{P3} Rendering stack}
\begin{itemize}[leftmargin=1.8em, itemsep=2pt]
  \sq Install Dr.Jit / Mitsuba~3 for reference renders
  \sq Install Blender headless (\texttt{bpy}) for dataset generation
  \sq Smoke test: render Cornell box with Blender Cycles (512\,spp, HDR output)
\end{itemize}

\subsection{\role{P4} Repository and paper}
\begin{itemize}[leftmargin=1.8em, itemsep=2pt]
  \sq Create public GitHub repo (suggested name: \texttt{sparserenderformer})
  \sq Add MIT or Apache~2.0 \texttt{LICENSE} file
  \sq Write README skeleton: installation, quick-start, citation block
  \sq Register account on target submission system (CMT / OpenReview --- confirm venue first)
\end{itemize}

%-----------------------------------------------------------------------
\section{Phase~2\quad Data Pipeline \textnormal{(Start Week~2, run in background throughout)}}

\critical{Do not wait until implementation is done.} Launch renders on the GPU as soon as
Blender is installed and let them accumulate overnight throughout Phases~3 and~4.

\subsection{\role{P3} Dataset construction}
\begin{itemize}[leftmargin=1.8em, itemsep=2pt]
  \sq Get access to RenderFormer's Objaverse subset (see repo README for download link)
  \sq Filter Objaverse by triangle count; write remesh script for
      $N \in \{4\text{k},\,8\text{k},\,16\text{k},\,32\text{k}\}$
  \sq Set up Blender Cycles pipeline: 512\,spp, HDR output, consistent lighting
\end{itemize}

\subsection{\role{P1} Dataset generation}
\begin{itemize}[leftmargin=1.8em, itemsep=2pt]
  \sq Generate $\geq\!1{,}000$ training scenes at $N = 16\text{k}$ triangles
  \sq Generate 50--100 held-out evaluation scenes at $N = 32\text{k}$ triangles
  \sq Write data loader that handles variable triangle counts
\end{itemize}

%-----------------------------------------------------------------------
\section{Phase~3\quad Implementation \textnormal{(Weeks~2--7, P1 + P2)}}

\subsection{\role{P1} BVH and coarse branch}
\begin{itemize}[leftmargin=1.8em, itemsep=2pt]
  \sq Implement BVH construction (SAH splitting, block size $b = 32$)
  \sq Implement coarse branch: area-weighted average pooling within BVH leaf nodes
  \sq Extend 3D RoPE to coarse branch using block centroid as position
  \sq Unit test: coarse branch output shape and gradient flow
\end{itemize}

\subsection{\role{P2} Selected and local branches}
\begin{itemize}[leftmargin=1.8em, itemsep=2pt]
  \sq Implement selected branch: top-$k$ BVH block selection per query via dot-product scores
  \sq Implement local branch: 3D kNN ($k = 64$) per query
  \sq Add normal-coherent mask with STE: hard threshold $\mathbf{n}_i \cdot \mathbf{n}_j < -0.5$
      in forward pass; sigmoid soft mask ($s = 20$) in backward pass for gradient flow to normals
  \sq Unit test: selected branch retrieves expected blocks on a synthetic scene
  \sq Unit test: local branch kNN distances and mask correctness
\end{itemize}

\subsection{\role{P1 + P2} Integration}
\begin{itemize}[leftmargin=1.8em, itemsep=2pt]
  \sq Implement gated combination (Eq.~5): per-head learned scalars
      $\lambda_1, \lambda_2, \lambda_3$ with softmax normalisation
  \sq Implement sparse view-dependent cross-attention: frustum culling + kNN ($k_r = 256$)
  \sq Fork and adapt RenderFormer codebase to slot in the GP-Sparse module
  \sq \critical{Sanity check:} at $N = 4{,}096$ with $k_\text{sel} = 128$,
      GP-Sparse output must match dense attention exactly
\end{itemize}

\subsection{\role{P2} Optional (if time allows)}
\begin{itemize}[leftmargin=1.8em, itemsep=2pt]
  \sq Profile naive PyTorch masking first; port NSA Triton kernel from DeepSeek-V3
      repo only if this is the bottleneck
\end{itemize}

%-----------------------------------------------------------------------
\section{Phase~4\quad Experiments \textnormal{(Weeks~5--9, overlaps Phase~3 tail)}}

Start running experiments on completed branches as soon as each passes unit tests.
Do not wait for full integration.

\subsection{\role{P4} Baselines}
\begin{itemize}[leftmargin=1.8em, itemsep=2pt]
  \sq RenderFormer pretrained, $N = 4{,}096$ \quad PSNR = \underline{\hspace{2.5cm}}
  \sq RenderFormer fine-tuned, $N = 8{,}192$ (dense, max memory)
      \quad PSNR = \underline{\hspace{2.5cm}}
\end{itemize}

\subsection{\role{P1 + P2} Main results}
\begin{itemize}[leftmargin=1.8em, itemsep=2pt]
  \sq Ours $N = 8{,}192$ \quad
      PSNR = \underline{\phantom{XX}}\,dB \quad
      Memory = \underline{\phantom{XX}}\,GB \quad
      Time = \underline{\phantom{XX}}\,ms
  \sq Ours $N = 16{,}384$ \quad
      PSNR = \underline{\phantom{XX}}\,dB \quad
      Memory = \underline{\phantom{XX}}\,GB \quad
      Time = \underline{\phantom{XX}}\,ms
  \sq Ours $N = 32{,}768$ \quad
      PSNR = \underline{\phantom{XX}}\,dB \quad
      Memory = \underline{\phantom{XX}}\,GB \quad
      Time = \underline{\phantom{XX}}\,ms
  \sq Log-log memory vs.\ $N$ plot (RF dense vs.\ GP-Sparse)
  \sq Log-log render time vs.\ $N$ plot (RF dense vs.\ GP-Sparse)
\end{itemize}

\subsection{\role{P4} Ablations}
\begin{itemize}[leftmargin=1.8em, itemsep=2pt]
  \sq Remove coarse branch, $N = 16\text{k}$ \quad PSNR = \underline{\hspace{2.5cm}}
  \sq Remove selected branch, $N = 16\text{k}$ \quad PSNR = \underline{\hspace{2.5cm}}
  \sq Remove local branch, $N = 16\text{k}$ \quad PSNR = \underline{\hspace{2.5cm}}
  \sq Block size $b \in \{16, 32, 64\}$ at $N = 16\text{k}$
  \sq $k_\text{sel} \in \{2, 4, 8\}$ at $N = 16\text{k}$
\end{itemize}

\subsection{\role{P3} Visual results}
\begin{itemize}[leftmargin=1.8em, itemsep=2pt]
  \sq Render visual gallery: Cornell box, 3 Objaverse scenes, 1 room scene
  \sq Side-by-side comparison figures: RF dense vs.\ GP-Sparse at $N = 32\text{k}$
\end{itemize}

%-----------------------------------------------------------------------
\section{Phase~5\quad Paper Completion \textnormal{(Weeks~7--10, overlaps experiments)}}

Start writing method, related work, and setup sections during implementation.
Fill in numbers as they arrive --- do not wait for every experiment to finish.

\subsection{\role{P4} Fill \texttt{\textbackslash{}todo} items in \texttt{overleaf/main.tex}}
\begin{itemize}[leftmargin=1.8em, itemsep=2pt]
  \sq Abstract: PSNR gain (dB), memory reduction (\%), render time reduction (\%)
  \sq Table~1: main results --- all PSNR / SSIM / LPIPS numbers
  \sq Table~2: branch ablation --- all 7 subsets at $N = 16{,}384$
  \sq Figure: memory/time scaling log-log plot
  \sq \S4.3: complexity ratio (XX$\times$ fewer token pairs)
  \sq \S5.3: fine-tuning epoch count
  \sq Conclusion: update all numbers
\end{itemize}

\subsection{\role{P4} Writing pass (CHECKLIST \S8b rules, in order)}
\begin{itemize}[leftmargin=1.8em, itemsep=2pt]
  \sq Abstract: no \texttt{\textbackslash{}cite\{todo\}} remain; target 150--220 words
  \sq Introduction: no consecutive sentences starting with the same word
  \sq Related Work: same sentence-start check
  \sq Method (\S4): all equations referenced in text
  \sq Experiments (\S5): no vague descriptors; report $\pm$std where applicable
  \sq Conclusion: no new results introduced; $\leq 150$ words
  \sq Dash audit:
      \texttt{grep -n "{-}{-}{-}" overleaf/main.tex | grep -v "\^{}\%"} must return zero hits
  \sq Banned-word grep: \emph{remarkable, impressive, strong, novel, intricate,
      delve, underscore, pivotal}
\end{itemize}

\subsection{\role{All} Remaining sections}
\begin{itemize}[leftmargin=1.8em, itemsep=2pt]
  \sq Acknowledgments: funding source, compute credits, dataset access
  \sq Ethics / Broader Impact statement
  \sq Supplementary material: extended ablations, additional renders, implementation details
\end{itemize}

%-----------------------------------------------------------------------
\section{Phase~6\quad Submission \textnormal{(Week~11)}}

\begin{itemize}[leftmargin=1.8em, itemsep=2pt]
  \sq Confirm venue and hard deadline: \underline{\hspace{5cm}}\\
      (SIGGRAPH Asia~2026, 3DV~2026, or ICCV~2027)
  \sq Submit arXiv preprint at least 1~week before conference deadline
  \sq Switch to anonymous mode in \texttt{main.tex}: toggle
      \texttt{ConferencePaper} $\to$ \texttt{ConferenceSubmission}
      and remove the author block
  \sq Page limit check: 8~pages + references
  \sq Upload fine-tuned checkpoint to HuggingFace Hub
  \sq Tag \texttt{v1.0} GitHub release at submission time
\end{itemize}

\noindent\textbf{After arXiv is live:}
\begin{itemize}[leftmargin=1.8em, itemsep=2pt]
  \sq Post to r/MachineLearning and r/computergraphics
  \sq Email RenderFormer authors (Chong Zeng, Microsoft Research Asia) with a
      one-paragraph note and preprint link; they may share feedback or repost
  \sq Add paper to Papers With Code
\end{itemize}

%-----------------------------------------------------------------------
\section{Compute Credits (apply as early as possible)}

\begin{itemize}[leftmargin=1.8em, itemsep=2pt]
  \item AWS Research Credits ---
        \url{https://aws.amazon.com/research-credits}
  \item Google TPU Research Cloud ---
        \url{https://sites.research.google/trc}
  \item Microsoft AFMR (Accelerate Foundation Models Research):\\
        {\footnotesize\url{https://microsoft.com/en-us/research/collaboration/accelerate-foundation-models-research}}
\end{itemize}

%-----------------------------------------------------------------------
\section{Optimised Timeline}

\begin{center}
\renewcommand{\arraystretch}{1.3}
\begin{tabular}{@{}clll@{}}
  \toprule
  \textbf{Weeks} & \textbf{Phase} & \textbf{Who} & \textbf{Note} \\
  \midrule
  1     & Learning sprint           & All 4             & Compressed to 1 focused week \\
  1--2  & Engineering setup         & P1 P2 P3 P4       & Runs in parallel with tail of Phase~0 \\
  2+    & Data renders (background) & P1 + P3           & GPU runs overnight throughout \\
  2--7  & Implementation            & P1 + P2           & 6 weeks; realistic for beginners \\
  5--9  & Experiments               & P1 P2 run;        & Start per branch as it passes tests \\
        &                           & P3 P4 visualise   & \\
  7--10 & Paper writing             & P4 leads          & Write while experiments run \\
  11    & Submission                & All               & \\
  \bottomrule
\end{tabular}
\end{center}

\vspace{6pt}
\noindent\textbf{Why 11 weeks instead of 13:}
\begin{itemize}[leftmargin=1.8em, itemsep=2pt]
  \item Learning compressed from 2~weeks to 1~week (save 1~week)
  \item Data renders overlap with implementation instead of following it (save 1~week)
  \item Paper writing overlaps with experiments instead of following them (save 1~week)
\end{itemize}

\end{document}
